% A.44: full independent parent proof/source audit and exact code rerun passed.
\subsection{Dihedral covers of two-power degree}

\begin{statement}{A.44}{Proved}{The four-cusp two-power dihedral family at every coefficient prime}
Let \(m=2^s\), \(s\ge1\), and put
\[
 B_m:\ v^2=w^{2m}-1,
 \qquad Y_m=B_m\setminus\{A,B,C,D\},               \tag{A.44.1}
\]
where \(A,B\) have \(w=0\), \(v=i,-i\), and
\(C,D\) are the two infinity points with
\(v/w^m\to1,-1\), respectively. Then \(B_m\)
has genus \(m-1\), and \(Y_m\to U\) is a finite
etale geometric dihedral cover of degree \(2m\).
Its proper Jacobian has potentially good reduction at
two, with special \(2\)-rank zero. Every difference
of two of the four cusps is killed by \(4m\).
Consequently the four cusp vertices of its pointed
stable reduction are joined entirely through rational
components.

For every coefficient prime \(\ell\), every cusp
packet at any one of these four cusps, and every
indicated open derived subgroup,
\[
                   \left(\pi_1(Y_{m,\bar K})^{(\ell)}\right)^{s(\Gamma_0)}
                                      =I_{P,\ell}.   \tag{A.44.2}
\]
This includes all induced original packets and permitted
open \(\Delta_0\) of Part A. It proves the fixed-group
part of Task A.5 for this second family of unbounded
genus. No normalizer classification for these new
higher-genus curves is asserted.
\end{statement}

\begin{proof}
\paragraph{The actual cover and its dihedral group.}
There is a map to \(E:y^2=x^3-x\) given by
\[
                 x=w^m,\qquad y=w^{m/2}v.            \tag{A.44.3}
\]
It has degree \(m\): the extension over \(K(x)\)
through \(K(w,v)\) has degree \(2m\), and the
intermediate elliptic function field has degree two.
It is the normalization of \(w^m=x\). The zeros
and poles of \(x\) on \(E\) have order two, so
there are two points above each of \(T=(0,0)\) and
\(O\), each with ramification index \(m/2\), and
no other ramification. Composing with the degree-two
isogeny \([1+i]\) in (A.43.8) gives the asserted
finite etale map to \(U\), of degree \(2m\),
with precisely the four displayed cusps removed.
The hyperelliptic polynomial has degree \(2m\) and
simple roots; hence the proper genus is \(m-1\).

Choose a primitive \(2m\)-th root \(\eta\) with
\(\eta^{m/2}=i\), after a finite extension, and set
\(\zeta_m=\eta^2\). The cyclic generator and a
lift of translation by \(T\) are
\[
 r(w,v)=(\zeta_m w,-v),
 \qquad h(w,v)=(\eta/w,-i v/w^m).                   \tag{A.44.4}
\]
The relation \(\eta^m=-1\) verifies the curve
equation, \(h^2=1\), and \(hrh=r^{-1}\).
Their induced maps on \(E\) are respectively the
identity and \((x,y)\mapsto(-x^{-1},yx^{-2})\).
They form a faithful group of order \(2m\) preserving
the composite map. Therefore the geometric Galois group
is dihedral of that order. For \(m=2\), this means
the group of order four. In every case it is a
\(2\)-group.

\paragraph{Explicit two-primary cusp divisors.}
The following divisors are computed on the smooth
compactification:
\[
 \begin{split}
 \operatorname{div}(w)&=A+B-C-D,\\
 \operatorname{div}\left(\frac{v-i}{v+i}\right)&=2m(A-B),\\
 \operatorname{div}(v-w^m)&=m(C-D).
 \end{split}                                       \tag{A.44.5}
\]
For the second identity use
\((v-i)(v+i)=w^{2m}\): at \(A\) one factor
is a unit and the other has order \(2m\), with
the opposite assertion at \(B\). Their poles at
the two infinities cancel in the quotient. For the
third use \((v-w^m)(v+w^m)=-1\); the respective
orders at \(C,D\) are \(m,-m\).
The first identity is the ordinary degree-two
hyperelliptic map at zero and infinity.

Twice \(m\) times the first divisor, plus the
second and minus twice the third, is \(4m(A-C)\).
With the second subtracted instead it is \(4m(B-C)\).
The third already kills \(C-D\) after multiplication
by \(m\). Thus every cusp difference is killed
by \(4m\), which is a power of two. We only need
this annihilator, and make no claim that all these
orders equal \(4m\).

\paragraph{Complex multiplication and good reduction.}
Let
\(\tau(w,v)=(\zeta_{2m}w,v)\). A holomorphic
basis is
\(w^j\,dw/v\), \(0\le j\le m-2\), so the
characters on this basis have exponents
\(1,\ldots,m-1\). On the full complex homology
we obtain each exponent
\[
              1,\ldots,m-1,m+1,\ldots,2m-1
\]
once. Thus the faithful endomorphism algebra generated
by \(\tau\) is
\[
 \mathbf Q[\tau]
   =\mathbf Q[t]\bigg/\frac{t^{2m}-1}{t^2-1}
   =\prod_{a=2}^{s+1} F_a,
                  \qquad F_a=\mathbf Q(\zeta_{2^a}).  \tag{A.44.6}
\]
Its dimension is \(2m-2=2g(B_m)\). The rational
idempotents split the Jacobian, up to isogeny after
a finite field extension, into abelian varieties
\(J_a\) with
\[
       \dim J_a=2^{a-2},\qquad
       F_a\hookrightarrow\operatorname{End}^0(J_a),
       \qquad [F_a:\mathbf Q]=2\dim J_a.            \tag{A.44.7}
\]
Indeed the \(F_a\)-summand in rational homology
has precisely that field's degree. No simplicity
of \(J_a\) is needed. These are abelian varieties
with complex multiplication in the usual degree-
\(2\dim J_a\) sense.

By \cite[Theorem 6(a), printed p.503]{fghSerreTate1968},
every \(J_a\) has potentially good reduction at two.
After a common finite extension so does their product
and its isogenous Jacobian. Good reduction is invariant
under isogeny by \cite[Corollary 2 to Theorem 1, printed p.493]{fghSerreTate1968};
equivalently apply the Neron--Ogg--Shafarevich criterion
to the rational Tate module at an auxiliary prime different from two. Make all the \(F_a\)-actions
defined before taking these good models. The same primary
proof of Theorem 6 specializes these field actions to
their good reductions using extension of endomorphisms
to the abelian schemes.

\paragraph{An elementary proof of special two-rank zero.}
For every \(a\ge2\), the polynomial
\[
                  \Phi_{2^a}(1+z)=(1+z)^{2^{a-1}}+1
                                                               \tag{A.44.8}
\]
is Eisenstein at two: the constant term is two, all
other nonleading coefficients are even, and the leading
coefficient is one. Hence
\(F_a\otimes\mathbf Q_2\) is a field of degree
\(2\dim J_a\).

Let \(\widetilde J_a\) be its good special fibre,
of dimension \(g_a\), and let
\[
 T_2^{\rm et}(\widetilde J_a)
       =\varprojlim_n\widetilde J_a[2^n](\bar\kappa),
 \qquad V_2^{\rm et}=T_2^{\rm et}\otimes\mathbf Q_2.
\]
Its dimension is the special \(2\)-rank \(f_a\),
which satisfies \(0\le f_a\le g_a\). One way
to see the upper bound is to factor \([2]\) through
relative Frobenius of purely inseparable degree
\(2^{g_a}\); its remaining Verschiebung factor
has degree \(2^{g_a}\), so the geometric kernel
of \([2]\) has at most \(2^{g_a}\) points.

The specialized \(F_a\)-action makes
\(V_2^{\rm et}\) a unital module over
\(F_a\otimes\mathbf Q_2\). If it were nonzero,
its \(\mathbf Q_2\)-dimension would be at least
\([F_a:\mathbf Q]=2g_a\), contradicting
\(f_a\le g_a\). Therefore \(f_a=0\).
Two-rank is additive in products and invariant under
isogeny, so the good special fibre of the full Jacobian
has two-rank zero. This argument requires neither a
Shimura--Taniyama slope computation nor a claim about
its complete Newton polygon.

\paragraph{The marked paths and the actual fixed groups.}
Take a semistable model of \(B_m\) after this finite
extension. Its dual graph is a tree, because the torus
rank of its semistable Jacobian is the graph's first
Betti number, whereas its Jacobian has good reduction.
This is the same compact-type criterion used in A.25
and A.40: the generalized special Jacobian is an
extension of the product of the component Jacobians
by \(\mathbf G_m^{b_1}\), as follows by gluing
line bundles along the nodes.

All cusp differences in (A.44.5) are two-primary
torsion. The special Jacobian has no nonzero geometric
two-primary torsion, so all four cusp Abel images
specialize to the same point. The extended Abel map
and its componentwise description in A.41 now show
that the entire path between every pair of marked
cusps consists of rational components. In particular
for each selected cusp there is a distinct cusp
satisfying the hypotheses of A.41--A.42. They prove
(A.44.2) at every coefficient prime.

At the coefficient prime two there is also an
independent direct verification: the geometric Galois
closure of this cover is the displayed dihedral
\(2\)-group. Thus A.22 embeds its full geometric
pro-two group as an open subgroup of the original
one and gives exactly the same inertia fixed group.
For odd coefficient primes the prime-to-residue-
characteristic statement A.41 already suffices.
The field and packet restrictions are precisely
those handled in A.41.

The exact program \texttt{code/a\_two\_power\_dihedral.py}
checks the curve and deck identities, the cusp-divisor
integer combinations, the multiplicity-one character
polynomial, and the Eisenstein certificates for several
powers of two. Their proofs above are uniform in \(s\).
\end{proof}
